% !TEX root = ../main.tex
% =====================================================================
% CAPITOLO 8 — CONCLUSIONI E SVILUPPI FUTURI
% =====================================================================

\chapter{Conclusioni e sviluppi futuri}
\label{chap:conclusioni}

% =====================================================================
\section{Sintesi dei contributi}
\label{sec:concl-sintesi}
% =====================================================================

Questa tesi ha affrontato il problema dell'adozione enterprise dell'AI generativa attraverso la progettazione, 
l'implementazione e la valutazione di un sistema multi-agente orchestrato, 
realizzato per Coop Alleanza 3.0. 

I contributi sono tre, in ordine di centralità.

\paragraph{1. Suite di agenti verticali specializzati.} I Capitoli~\ref{chap:agenti-architettura} e~\ref{chap:agenti-implementazione} hanno descritto la progettazione e l'implementazione di quattro agenti,
ciascuno con la propria forma di grounding (Web, RAG, Text-to-SQL), serviti da un'infrastruttura cloud-native interamente descritta in Terraform, 
con autenticazione federata, catena di sicurezza pre e post generazione e pipeline CI/CD automatizzata.

La scelta di agenti verticali indipendenti ha questi punti di forza: ogni agente può essere aggiornato, scalato e monitorato indipendentemente; 
il prompt di ciascun agente è ottimizzato per il proprio dominio; è dimostrata, con la coesistenza di Gemini e Claude in uno stesso agente,
la possibilità di evitare il vendor lock-in.

\paragraph{2. Orchestratore Orion.} I Capitoli~\ref{chap:orion-architettura} e~\ref{chap:orion-implementazione} hanno presentato Orion, 
un orchestratore LLM-based che instrada le richieste dell'utente verso l'agente specializzato più appropriato. 
L'architettura si fonda su un routing step che produce un output strutturato con cinque campi (agente target, complessità, confidence, reasoning, clarification), 
eseguito da Gemini~2.5 Flash con guided generation. Il meccanismo di confidence con soglia configurabile consente al sistema di porre domande di chiarimento quando l'ambiguità è alta, 
in un'implementazione del principio di \emph{active learning} applicato al routing.

Il proxy normalizza le risposte eterogenee dei sub-agenti verso un formato uniforme per il frontend, 
e l'autenticazione service-to-service basata sull'ingress Cloud Run elimina la complessità del token forwarding. 
Il pannello di amministrazione consente di visualizzare i log di routing, modificare il system prompt del router senza rilascio di codice, 
e lanciare la valutazione automatica direttamente dal browser.

\paragraph{3. Metodologia di valutazione del routing.} Il Capitolo~\ref{chap:orion-valutazione} ha introdotto una metodologia per la valutazione sistematica del routing, 
basata su un gold standard di 120~test case, uno strumento di esecuzione automatica integrato nell'admin panel e cinque metriche dichiarate a priori
(agent accuracy, complexity accuracy, clarification rate, mean confidence sui casi corretti e mean confidence sui casi errati). 
La metodologia ha consentito di misurare l'effetto di un'iterazione del system prompt del router: 
agent accuracy stabile al $99{,}2\%$ e complexity accuracy dall'$88{,}3\%$ al $93{,}3\%$ tra le versioni v1.0 e v1.1, sopra i rispettivi target.

A complemento della valutazione del routing, il Capitolo~\ref{chap:orion-valutazione} riporta anche una valutazione quantitativa della qualità dei sub-agenti: 
execution accuracy dell'$80$--$84\%$ per gli agenti Text-to-SQL e, per il RAG, faithfulness $0{,}73$ con answer rate del $100\%$ sulle domande in corpus (Sez.~\ref{sec:eval-agenti}), 
e una valutazione economica indicativa (Sez.~\ref{sec:eval-costi}) che dà sostanza alla motivazione di costo: nello scenario illustrativo considerato, 
una piattaforma a consumo risulta di circa un fattore 40 più economica delle licenze SaaS nominali per l'intera popolazione, 
vantaggio che deriva dal modello di tariffazione (uso effettivo contro disponibilità per utente) più che dalla tecnologia.

% =====================================================================
\section{Generalizzabilità del pattern}
\label{sec:concl-generalizzabilita}
% =====================================================================

L'architettura proposta non è legata in modo essenziale al dominio dell'organizzazione. I pattern emersi (agenti verticali con grounding eterogeneo, orchestratore LLM-based con confidence calibration, 
valutazione sistematica del prompt) sono riutilizzabili in ogni organizzazione con caratteristiche di questo tipo:

\begin{itemize}
  \item Infrastruttura cloud consolidata (non necessariamente GCP: i pattern architetturali sono portabili su Azure o AWS) o specifica per AI on premises.
  \item Difficoltà economica o normativa nell'adottare licenze SaaS nominali per l'intera popolazione aziendale.
  \item Patrimonio informativo eterogeneo: documenti non strutturati (policy, procedure), dati strutturati (ticket, sondaggi) e necessità di accesso a informazioni pubbliche attuali.
  \item Vincoli di compliance e sovranità del dato che richiedono il controllo completo dell'infrastruttura e del flusso dei dati.
\end{itemize}

Il codice sorgente completo del progetto, anonimizzato e reso indipendente dal contesto organizzativo specifico,
è pubblicato nel repository \url{https://github.com/dpiccinelli84/unibo-thesis-src} per consentire la riproducibilità. 
Il repository include anche il modulo \texttt{evaluation/}, con gli script e i gold standard delle valutazioni di qualità degli agenti (Sez.~\ref{sec:eval-agenti}) 
e il calcolatore del modello di costo (Sez.~\ref{sec:eval-costi}).

% =====================================================================
\section{Sviluppi futuri}
\label{sec:concl-futuri}
% =====================================================================

\subsection{Da router LLM-based a router appreso}
\label{subsec:concl-routellm}

La direzione più naturale è la transizione dal router LLM-based attuale a un router \emph{appreso}~\cite{ong2024routellm}, 
addestrato su dati di preferenza derivati dal traffico reale. 
Il log di routing di Orion, che registra ogni decisione con prompt, agente selezionato, confidence e reasoning, 
è il dataset di training naturale per un classificatore supervisionato. 
Il vantaggio atteso è duplice: riduzione della latenza (inferenza locale vs chiamata API) e riduzione del costo (nessuna chiamata Flash per il routing). 
Il prerequisito è un volume di traffico sufficiente a produrre un dataset bilanciato su tutti gli agenti, condizione non ancora soddisfatta nella fase attuale del progetto.

\subsection{Estensione della suite a nuovi domini}
\label{subsec:concl-nuovi-domini}

L'architettura modulare consente l'aggiunta di nuovi agenti specializzati senza modifiche strutturali: ogni nuovo agente richiede un backend, un frontend, 
l'aggiornamento del prompt del router (per includere la descrizione del nuovo agente) e una regola nel bilanciatore di carico. 
Domini candidati includono l'analisi dei turni di lavoro (già in fase prototipale), la gestione documentale SharePoint e l'analisi dei dati di vendita.

\subsection{Valutazione end-to-end su traffico reale}
\label{subsec:concl-valutazione-reale}

La valutazione del routing (Sez.~\ref{sec:eval-risultati}) opera in isolamento, mentre quella dei sub-agenti (Sez.~\ref{sec:eval-agenti}) è già end-to-end ma su un gold standard statico. 
Un'estensione naturale è la valutazione \emph{continua} sul traffico reale, con revisione umana periodica delle decisioni di routing e delle risposte degli agenti. 
Il framework AgentBench~\cite{liu2023agentbench} fornisce un modello metodologico per questo tipo di valutazione; l'infrastruttura di logging già in essere (routing logs, chat logs, feedback utente) 
fornisce i dati necessari.

% =====================================================================
\section{Considerazione finale}
\label{sec:concl-finale}
% =====================================================================

L'adozione enterprise dell'AI generativa non è la scelta binaria tra ``comprare una licenza SaaS'' e ``non adottarla'': 
è un problema di \emph{architettura dei sistemi informativi} a tutti gli effetti. 
Posto in questi termini, ammette risposte ingegnerizzate che sono economicamente sostenibili, 
normativamente difendibili e tecnicamente controllabili. 
Questa tesi ha mostrato una di queste risposte e ha fornito gli strumenti per valutarla in modo sistematico. 
Il lavoro non pretende di aver risolto il problema in modo definitivo, ma di aver dimostrato che l'approccio è praticabile, misurabile e, 
per il contesto specifico di Coop Alleanza, preferibile alle alternative considerate.
